\documentclass[11pt,a4paper]{article}
\usepackage[T1]{fontenc}
\usepackage[utf8]{inputenc}
\usepackage{lmodern,microtype}
\usepackage[margin=27mm]{geometry}
\usepackage{amsmath,amssymb,amsthm,nicefrac}
\usepackage[hidelinks]{hyperref}
\usepackage[numbers]{natbib}
\hypersetup{pdftitle={Justified Upper Quota and Adams-Thiele},pdfauthor={Oliviero Nardi}}
\setlength{\parindent}{0pt}
\setlength{\parskip}{5pt}
% The notation of the Preliminaries and Representation chapters.
\newcommand{\voters}{N}
\newcommand{\candidates}{C}
\newcommand{\approfile}{A}
\newcommand{\app}[1]{A(#1)}
\newcommand{\profiles}{\mathcal A}
\newcommand{\committee}{W}
\newcommand{\committees}{\mathcal W}
\newcommand{\kcommittees}{\mathcal W_k}
\newcommand{\weight}{\omega}
\newcommand{\num}{\ell}
\newcommand{\nats}{\mathbb N}
\newcommand{\nonNegReals}{\mathbb R_{\ge0}}
\newcommand{\Thiele}{\textnormal{\textsc{Thl}}}
\newcommand{\score}{\operatorname{score}}
\newcommand{\quota}{q}
\newcommand{\Adams}{\textnormal{Adams-Thiele}}
\newcommand{\vect}[1]{\left(#1\right)}
\newcommand{\ceil}[1]{\left\lceil#1\right\rceil}
\newcommand{\powerset}[1]{2^{#1}}
\newcommand{\set}[1]{\{#1\}}
\newcommand{\setcompr}[2]{\{#1:#2\}}
\DeclareMathOperator*{\argmax}{arg\,max}
\newtheorem{definition}{Definition}
\newtheorem*{adamsjuq}{Theorem (\texttt{AdamsJUQ})}
\newtheorem*{onlyadamsjuq}{Theorem (\texttt{OnlyAdamsJUQ})}
\title{Justified Upper Quota and Adams-Thiele\\[3pt]
  \large Definitions and theorem statements}
\author{Oliviero Nardi}
\date{November 2026}
\begin{document}
\maketitle
\vspace{-1em}
This note states \texttt{AdamsJUQ} and \texttt{OnlyAdamsJUQ} from the
accompanying Lean interface.
The original paper is my work with Martin Lackner, \emph{Avoiding
Overrepresentation: Upper Quota Axioms for Committee Voting}, presented at
M-PREF 2025~\cite{lacknernardi2025}. This work was developed further in our
EC 2026 paper~\cite{baychkov2026} and my thesis~\cite{nardi2026}.
The definitions below use the notation and wording of the Preliminaries and
Representation chapters of the thesis. More details, examples, and proofs
are available in these sources.
I wrote the Lean definitions and theorem statements; the proofs were
autoformalised with Codex and checked by Lean.
Codex was given the proofs included in my thesis as a starting point,
but the formalised proofs may differ from the arguments in the thesis.

\section{Context}

These results pertain to the theory of multi-winner elections.
Here, we are given a set of candidates and must select a fixed number of them.
For example, consider parliamentary or committee elections. Voters express their preferences
via approval ballots; that is, they can approve as many candidates as they wish. 
Our contribution is a formalisation of the notion of \emph{overrepresentation}, that is,
a scenario in which an interest group has a disproportionally-large influence on the result of the election.
For example, this models situations in which a minority ends up controlling a majority of the parliament.
Here, we formalise a theorem stating that, among a class of natural voting rules,
there is exactly one that avoids overrepresentation (under our definition).

\section{The Model and the Rules}
The set $\nats=\{1,2,\ldots\}$ denotes the positive integers, and let
$[\num]=\{1,\ldots,\num\}$ for $\num\in\nats$.
We indicate as $\candidates=\{a,b,\ldots\}$ the finite set of
\emph{candidates} with $|\candidates|=m$. Furthermore, we assume that there
are $n\in\nats$ \emph{voters}, and that the set of voters is $\voters=[n]$.
We are also given a target \emph{committee size}, indicated as $k\in[m]$.
The preferences of the voters are given by the so-called \emph{profile}
$\approfile\colon\voters\to\powerset{\candidates}$. In other words, each
voter $i\in\voters$ submits an approval ballot $\approfile(i)\subseteq\candidates$.
We typically denote
$\approfile$ with a vector notation, that is,
$\approfile=\vect{\app{1},\ldots,\app{n}}$.
Given a candidate $c\in\candidates$, we define
$\voters(c)=\{i\in\voters:c\in\app{i}\}$ as the set of its supporters.
Write $\profiles$ for the set of all profiles.

A \emph{committee} is a set $\committee\subseteq\candidates$ of size at most
$k$. If $|\committee|=k$, we say that $\committee$ is an \emph{exhaustive}
committee. We denote as $\committees$ the set of all committees and as
$\kcommittees$ the set of all exhaustive committees. Given that many profiles
naturally lead to ties, we formalise a voting rule as a function $F$ that returns
a nonempty set of committees, the set of (tied) winning committees:
\[F:\profiles\rightarrow\powerset{\kcommittees}\setminus\{\varnothing\}.\]
Given two voting rules $F$ and $G$, we say that $F$ refines $G$ if, for every
$\approfile\in\profiles$, it holds that $F(\approfile)\subseteq G(\approfile)$.

We now define our family of voting rules. First, we need some auxiliary notions:

\begin{definition}\label{def:weight-functions}
	A \emph{weight function} is a function
  $\weight\colon\nats\rightarrow\nonNegReals$ such that $\weight(1)=1$ and, for
  every $\num\in\nats$, $\weight(\num)\geq \weight(\num+1)$. We sometimes write
  weight functions as vectors.
\end{definition}

The following is a standard family of voting rules \cite{abcbook}, called Thiele rules.

\begin{definition}\label{def:Thiele}
  Given $\approfile\in\profiles$, $i\in\voters$ and
  $\committee\in\kcommittees$, we define
  \[
    \score_{\weight}(\app{i}, \committee) = \sum_{\num=1}^{\lvert\app{i}\cap
    \committee\rvert}\weight(\num)\quad\quad\text{and}\quad\quad\score_{\weight}(\approfile,
    \committee)=\sum_{j\in \voters}\score_{\weight}(\app{j}, \committee).
  \]
  The \emph{Thiele rule} corresponding
  to $\weight$ is then
  \[
    \Thiele_\weight(\approfile) =
    \argmax_{\committee\in\kcommittees}\Big(\score_{\weight}(\approfile,\committee)\Big).
  \]
\end{definition}

We extend this class to \emph{composite} Thiele rules.
Intuitively, composite Thiele rules are defined through a sequence of Thiele
rules and work in stages. First, we compute the result of the first Thiele rule.
Then, at each subsequent stage, we apply the next rule to break the ties of
the previous stage.

\begin{definition}
  Recall the definition of weight functions
  and Thiele rules (Definitions~\ref{def:weight-functions} and~\ref{def:Thiele}).
  Recall that the Thiele rule $\Thiele_\weight$ picks
  all committees $\committee$ maximising
  \[
\score_{\weight}(\approfile,
    \committee)=
    \sum_{i\in\voters}\sum_{\num=1}^{\lvert\app{i}\cap
    \committee\rvert}\weight(\num).
  \]
  Given two weight functions $\weight_1$ and $\weight_2$, we define the composition
  \[
    \left(\Thiele_{\weight_2}\circ \Thiele_{\weight_1}\right)(\approfile) = \argmax_{\committee\in
    \Thiele_{\weight_1}\vect{\approfile}}\ \score_{\weight_2}\mathclose{\vect{\approfile, \committee}}.
  \]
  For any sequence of weight functions $\boldsymbol{\weight}=\vect{\weight_1,\weight_2,\ldots}$
  of length $\num\in\nats\cup\{\infty\}$,
  the
  corresponding \emph{$\num$-composite Thiele rule} is
  \[
    \Thiele_{\boldsymbol{\weight}}(\approfile) = \bigcap_{i\in[\num]}
    (\Thiele_{\weight_i}\circ \cdots \circ \Thiele_{\weight_1})(\approfile).
  \]
  Note that if $\boldsymbol{\weight}=\vect{\weight_1,\ldots,
  \weight_\num}$ is finite then the
  above collapses to
  \[
    \Thiele_{\boldsymbol{\weight}}(\approfile) =
    \left(\Thiele_{\weight_\num}\circ \cdots \circ \Thiele_{\weight_1}\right)(\approfile)
  \]
  and that this is well-defined even for sequences of length $\num = \infty$ as there are
  only finitely many committees.
\end{definition}

We finally introduce our main voting rule.

\begin{definition}\label{def:adams-thiele}
	\emph{\Adams} is the $2$-composite Thiele rule given by
  $\weight_1=\vect{1,0,0,\ldots}$ and
  $\weight_2=\vect{1,1,\nicefrac{1}{2},\,\nicefrac{1}{3},
  \,\ldots}$.
\end{definition}
Observe that \Adams{} first maximises the number of covered voters, that is,
voters approving at least one candidate in the committee (just as the Chamberlin-Courant rule does \cite{abcbook}).
Then, it proceeds with a weight function similar to Proportional Approval Voting rule \cite{abcbook}.

\section{Justified Upper Quota and the Formalised Theorems}
For a set of voters $\voters^\prime\subseteq\voters$, we define its quota as
\[
  \quota(\voters^\prime)
  =\frac{k\cdot|\voters^\prime|}{|\voters|}.
\]
We call $\ceil{\quota(\voters^\prime)}$ the upper quota of
$\voters^\prime$. We finally define $\quota_c=\quota(\voters(c))$ for a
candidate $c\in\candidates$.

\begin{definition}\label{def:JUQ}
	Committee $\committee$ satisfies \emph{JUQ} if,
  whenever there is a $c\in \committee$ such that
  \begin{align*}
    \lvert\app{i}\cap\committee\rvert
    &> \ceil{\quota_c}
    &&\text{for all } i\in\voters(c), \\
    \intertext{then there exists no $d\in \candidates\setminus \committee$
    and nonempty $\voters^\prime\subseteq\voters(d)$ such that}
    \lvert\app{i}\cap((\committee\cup\{d\})\setminus\{c\})\rvert
    &\le \ceil{\quota(\voters^\prime)}
    &&\text{for all } i\in\voters^\prime.
  \end{align*}
\end{definition}
Note that we allow for arbitrary subsets $\voters^\prime$ of $\voters(d)$ in
our definition: for us it is sufficient that there exists some nonempty
subgroup of the approvers of $d$ who could be awarded an extra candidate without
breaching their upper quota. A rule satisfies JUQ if every committee it returns
does so.

\begin{adamsjuq}
	\Adams{} satisfies JUQ for every choice of $m\in\nats$, $n\in\nats$, and $k\in[m]$.
\end{adamsjuq}

\begin{onlyadamsjuq}
Fix a sequence of weight functions
$\boldsymbol{\weight}$. If $\Thiele_{\boldsymbol{\weight}}$ satisfies JUQ
for every choice of $m\in\nats$, $n\in\nats$, and $k\in[m]$, then
$\Thiele_{\boldsymbol{\weight}}$ is a refinement of \Adams.
\end{onlyadamsjuq}

The same sequence $\boldsymbol{\weight}$ is fixed across all election sizes.

\bibliographystyle{plainnat}
\bibliography{references}
\end{document}
